\scanpage{178}\mspage{33}%
% archive sheet number: 178
% C original page filename: H134-06_134-2_scan178.tex
vérifier et d'expliciter comment les ``6 opérations'' habituelles sur les faisceaux (soit sur $X$, soit sur une réunion de strates de $X$) peuvent s'exprimer dans ce dictionnaire, dans le cas disons de strates non singulières (\unsure{autrement}, on aurait des difficultés avec les complexes dualisants, dont on aimerait qu'ils appartiennent à notre catégorie\unsure{.\,.\,.\,.\,.|\ldots}) \inserted{et de revérifier le formulaire connu de ces opérations}. Mais il apparaitra probablement que, pour bien pouvoir tout transcrire, il convient, plutôt \inserted{que} de regarder un diagramme du type
% Diagram physically within L11. A separate point at approximately (1229,905)
% lies right and below the right arrowhead: reproduced as a nontext mark,
% not as a terminal object or mathematical operator; its origin is uncertain.
\diagram{\begin{tikzcd}[ampersand replacement=\&,column sep=1.4em,row sep=1.2em]
\cdot \arrow[r] \arrow[d] \& {\raisebox{0pt}[0pt][0pt]{\parbox{0pt}{\figure{\begin{tikzpicture}[overlay]\fill (0.25,-0.18) circle[radius=0.25pt];\end{tikzpicture}}}}} \\
{} \& {}
\end{tikzcd}}
dans la catégorie homotopique formée avec la cat.\,. des ens. semi-simpliciaux, \inserted{de} regarder la catégorie des diagrammes d'ensembles $1/2$ simpliciaux, et à partir de là passer à une catégorie homotopique de
% Separate dark point at approximately (1157,1368), above-left of fractions.
% Its presence is clear; its origin and function are not assigned.
\raisebox{0pt}[0pt][0pt]{\parbox{0pt}{\figure{\begin{tikzpicture}[overlay]\fill (0,0.60) circle[radius=0.3pt];\end{tikzpicture}}}}fractions.

J'ai plus ou moins explicité dernièrement, en cogitant sur \deleted{\unsure{les|des|ces}}des fondements de ``topologie modérée'' (i.e. où on élimine d'emblée tous les phénomènes de sauvagerie), comment une stratification équisingulière%
% L21a: the smaller cancelled parenthesis is above équisingulière;
% its enclosing/attachment stroke turns down at the word's right end.
\inserted{\deleted{(\unsure{d'un} \unsure{espace} \unsure{modéré} \unsure{compact})}}, à strates non singulières disons, d'un ``espace modéré'' \inserted{compact}, donne naissance canoniquement à un diagramme \deleted{\unsure{fon|f}} d'espaces,
